%!TEX TS-program = xelatex
\documentclass[10pt,table,a4]{article}\usepackage[]{graphicx}\usepackage[]{color}


%%%%%%%%%%%%
% To insert a new page, copy in the following text and populate the space between \HEaderSingle and \PageFooter
% Populate the {brackets} with either the header title or footer text. 
% \HeaderDouble{FirstLineText}{Second Line Text} is also an option
% Use unique variable handles for text that must be replaced
% Highlight text that is "flexible" that will need to be subbed out with text from R
% Consider checking the GeneralTemplate for inspriration for tables etc. 

%	\section*{} % Introduction
%	\HeaderSingle{Company Profile}
%
%	This is the company report for Company.Name
%	\adjincludegraphics[width = .95\linewidth,trim={0cm 0cm 0cm .5cm},clip]{BubbleChart.png}
%	
%	Flex.Text.P01       % An example handle for text that needs to be written in R based on the results. 
%
%	\PageFooter{Company Profile}
%	\newpage
%%%%%%%%%%%%



\usepackage{alltt}
\usepackage{graphicx}
\usepackage{gensymb}
\usepackage[top=1cm, bottom=1.5cm, left=1.8cm, right=1.8cm]{geometry}
\usepackage[font=small]{caption}
\usepackage{adjustbox}
\usepackage{fancyhdr}
\usepackage{layout}
%\usepackage{booktabs}
%\usepackage{kpfonts}
\usepackage[explicit]{titlesec}
\usepackage{wrapfig}
\usepackage{tcolorbox}
\usepackage{xcolor}
\usepackage{setspace}
\usepackage{parskip}
\usepackage{tikz}
\usepackage{fontspec}
\usepackage{anyfontsize}
\usepackage{hyphenat}
\usepackage{multicol}
\usepackage{datetime}
\usepackage{fixltx2e}
\usepackage{mfirstuc}
\usepackage{tabularx}
\usepackage{microtype}
\usepackage{vwcol} 
\usepackage{parcolumns}
\usepackage{xcolor}
\usepackage{xpatch}
\usepackage[utf8]{inputenc}
\usepackage{hyperref}
\usepackage{enumitem} %allow changig space between \item
\usepackage{booktabs} % table 
\usepackage[sfdefault]{ClearSans}

% Colours
\definecolor{Blue1}{RGB}{47,73,110}
\definecolor{Yellow2}{RGB}{252, 190, 54}
\definecolor{Yellow3}{RGB}{254, 238, 207}

\definecolor{PearlBlack}{RGB}{8,12,18}
\definecolor{LightGray}{RGB}{208,208,208}

\definecolor{ColRed}{RGB}{244,123,115}
\definecolor{ColOrange}{RGB}{253,226,145}
\definecolor{ColYellow}{RGB}{255,255,204}
\definecolor{ColGreen}{RGB}{195,214,155}

\pagestyle{fancy}
\fancyhf{}
%\fancyhead[R]{\thepage}
\renewcommand{\headrulewidth}{0pt}

\setlength{\parskip}{20pt}
\righthyphenmin=2

\setmainfont{ClearSans}
\color{PearlBlack}



\newcolumntype{P}[1]{>{\centering\arraybackslash}p{#1}}

\newcommand{\nearr}[1]{$\color{#1}\nearrow$}
\newcommand{\searr}[1]{$\color{#1}\searrow$}
\newcommand{\uparr}[1]{$\color{#1}\uparrow$}

\setlength{\parskip}{5pt}

\hypersetup{
	colorlinks=true,
	linkcolor=cyan,
	filecolor=magenta,      
	urlcolor=blue,
}

\urlstyle{same}

\newcommand*{\PageHeadingSingleLine}{%
	\begin{tikzpicture}[remember picture,overlay]
	\node[anchor=north west,minimum width=.375cm,minimum height=1.2cm,fill=Blue1] (RB) at (-1.8,1.2){\Large };
	
	\end{tikzpicture}}
\newcommand*{\PageHeadingDoubleLine}{%
	\begin{tikzpicture}[remember picture, overlay]
	\node[anchor=north west,minimum width=.375cm,minimum height=1.9cm,fill=Blue1] (RB) at (-1.8,1.2){\Large };
	
	\end{tikzpicture}}
\newcommand{\HeaderSingle}[1]{
	\PageHeadingSingleLine 
	
	\vspace{-1.2cm}
	{\huge\textbf{\textcolor{Blue1}{#1}}}
	\vspace{.2cm}}
\newcommand{\HeaderDouble}[2]{
	\PageHeadingDoubleLine
	
	\vspace{-1.2cm}
	{\huge\textbf{\textcolor{Blue1}{#1 \\[2pt] #2}}}
	\vspace{.45cm}}
\newcommand*{\SectionHeadingBox}[1]{%
	\begin{tikzpicture}[remember picture, overlay]
	\node[anchor=north west,minimum width=.375cm,minimum height=#1,fill=Blue1] (RB) at (-1.2,-16){\Large };
	\end{tikzpicture}
	\vspace{.8cm}}
\newcommand{\SectionHeading}[2]{
	\SectionHeadingBox{3cm}
	
	\vspace{15.7cm}
	\textbf{{\Huge #1 \\[6pt]\Huge #2}}}
\newcommand{\SectionHeadingDouble}[3]{
	\SectionHeadingBox{4cm}
	
	\vspace{15.7cm}
	\textbf{{\Huge #1 \\[6pt]\Huge #2\\[6pt]\Huge #3}}}
\newcommand{\PageFooter}[1]{% The number indicates the sector..... 
	\begin{tikzpicture}[remember picture, overlay]
	\node[text=Blue1,above = .8cm, left = -9.43cm,font=\bf\small,align=center] at (current page.south){ #1 | \thepage};	
	\end{tikzpicture}
}

\newcommand{\KeyQuestions}{\textbf{\textcolor{Blue1}{\normalsize Key Questions}}}
\newcommand{\KeyQuestion}{\textbf{\textcolor{Blue1}{Key Question}}}
\newcommand{\Companyname}{\textbf{\textcolor{Blue1}{Company.Name}}}
\newcommand{\Subheader}[1] {\vspace{10pt}{\textbf{\textcolor{Blue1}{\large{#1}}}\vspace{5pt}}}
\newcommand{\Colortext}[1] {\vspace{0pt}{\textbf{\textcolor{Blue1}{\normalsize{#1}}}\vspace{3pt}}}
\newcommand{\ColortextLarge}[1] {\vspace{3pt}{\textbf{\textcolor{Blue1}{\normalsize{#1}}}\vspace{6pt}}}
\newcommand{\Fig}[1]{\textbf{\textcolor{Blue1}{Figure #1}}}
\newcommand{\FigCap}[1]{{\textcolor{Blue1}{#1}}}
\newcommand{\Tab}[1]{\textbf{\textcolor{Blue1}{Table #1}}}
\newcommand{\TabCap}[1]{{\textcolor{Blue1}{#1}}}
\renewcommand{\thefootnote}{\fnsymbol{footnote}}
\newcommand{\Footnotetext}[1] {\vspace{0pt}{\textcolor{Blue1}{\footnotesize{#1}}\vspace{3pt}}}

\setmainfont{ClearSans}
\color{PearlBlack}


% Box with a shaded border
\newtcolorbox{shadedbox}{colback=Blue1}

%%%%%%%%%% BEGIN %%%%%%%%%%
\begin{document}
	
	%%%%%%%%%% FRONT PAGE %%%%%%%%%%
	\section*{} % Front Page
	
	
	\thispagestyle{empty}
	\pagecolor{Blue1}
	
	\vspace{-2.6cm}
	
	\begin{tikzpicture}[remember picture, overlay]
	\vspace{7cm}
	\node[below=-5cm] (CN) at (current page.north){\adjincludegraphics[height=19cm,trim={0cm 0cm 0cm 0cm},clip]{CompanyReportGraphics/Cars.png}
	};
	
	\node[below = 5cm, right,align = left, white, text width=12cm](TN) at (5,-14.5){
		\Huge{\textbf{Climate  Scenario  Analysis}}\\[10pt]
		%		\vspace{2cm}
		%		{\baselineskip=20pt\Huge\raggedright{\textbf{InvestorName}\par}}
		\vspace{1.4cm}
		{\baselineskip=14pt\Huge\raggedright{\textbf{Company.Name}\par}}
		\vspace{1.4cm}
		{\baselineskip=12pt\Huge\raggedright{\textbf{Report Start.Year}\par}}};
	
	
	\end{tikzpicture}
	
	
	\begin{tikzpicture}[remember picture,overlay]
	\node[anchor=south west, yshift = 0cm ,minimum width=21.6cm,minimum height=4cm,fill=white] (RB) at (current page.south west){};
	\node[left= 7cm, above=1.cm] (WS) at (current page.south){\adjincludegraphics[height=1.8cm,trim={0cm 0cm 0cm 0cm},clip]{CompanyReportGraphics/Logo_front}};
	\node[ below = 2.15cm, minimum width = 3.6cm, minimum height =0.35cm, fill = Blue1] at (1.85,-25.25){};
	\end{tikzpicture}
	
	\vspace{24cm}
	\hspace{5cm}
	\begin{tikzpicture}
	%\draw (0, 0)
	\node[rounded corners = 5pt, fill = red, text = white, minimum height = 1cm]
	{\Large{\textbf{DRAFT -  DO NOT CIRCULATE OR USE TO INFORM DECISIONS}}};
	\end{tikzpicture}
	
	\newpage
	
	\pagecolor{white}
	
	%%%%%%%%%% 2 %%%%%%%%%%
	\section*{} % Legal Disclaimer
	
	\HeaderSingle{Legal Disclaimer}
	
	\small{	
		This Report has been prepared by the 2° Investing Initiative (2°ii) a leading not-for-profit think-tank on climate related metrics and policies in financial markets. The Report summarises our Company Climate Scenario Analysis (CCSA) in relation to \(\Companyname\)  (the Company). The CCSA is our limited 'point in time' estimate of the alignment between the Company's revealed business plans for its automotive production business in the period 2019-2024, versus the economic trends embodied in the International Energy Agency’s (IEA’s) 'World Energy Outlook' and 'Energy Technology Perspective' scenarios. This Report is made available through our PACTA portal and its use is subject to the Terms of Service agreed to by users of that portal. The methodology applied in the CCSA, its data inputs, assumptions and limitations, are set out in this Report and the Methodology Statement – Company Reports available at: \par www.transitionmonitor.com/company-reports/. 
		
		\textbf{Limitations and assumptions:}
		The CCSA does not purport to analyse all risks, opportunities or issues associated with climate change that may be relevant to the Company. Such issues may include (for example) physical or ecological impacts that may be caused by, or to, the assets and operations of the Company, and any climate-related litigation exposures.
		The CCSA utilises publicly-available information, and proprietary third party data obtained under licence, which 2°ii believes in good faith to be reliable. However, 2°ii makes no representation or warranty (express or implied) as to the completeness, accuracy or currency of such information or data, nor to the completeness, accuracy or currency of the information in this Report. 
		
		\textbf{No forecast or prediction:}
		The CCSA does not purport to generate, nor does this Report contain or comprise, forecasts or predictions. 2°ii neither makes nor implies any representation regarding the likelihood, risk or expectation of any future matter. In particular, 2°ii does not make any representation that by making the changes to the Company's revealed plans identified in the Report, the Company would be in compliance with any Sustainable Development Goals or the long-term warming goals of the Paris Agreement or achieve any stated climate targets. To the extent that any statements made or information contained in this Report might be considered forward-looking in nature, they are subject to risks, variables and uncertainties that could cause actual results to differ materially. You are cautioned not to place any reliance on any such forward-looking statements, which reflect our assumptions and methodology as applied to third party data and the Company's revealed business plans only as of the date of modelling or such earlier date as indicated in this Report. It is likely that the third party data, the Company's revealed plans, and the IEA scenarios will change in some way during the five-year time horizon, and our assumptions and methodology may also evolve and change during this time. 2°ii is not obliged to revise, or to publicly release any revisions to, this Report or to notify you if the data, revealed plans, assumptions or methodology change or become inaccurate.  
		
		\textbf{No financial advice:}
		The information contained in this Report is general in nature. It does not comprise, constitute or provide personal, specific or individual recommendations or advice, of any kind. In particular, it does not comprise, constitute or provide, nor should it be relied upon as, investment or financial advice, a credit rating, an advertisement, an invitation, a confirmation, an offer, a solicitation, an inducement or a recommendation, to buy or sell any security or other financial, credit or lending product, to engage in any investment strategy or activity, nor an offer of any financial service. \textbf{This Report does not purport to quantify, and 2°ii makes no representation in relation to, the performance, strategy, prospects, creditworthiness or risk associated with the Company or any investment therein, nor the achievability of any stated climate targets (of the Company, the defined market, an investor's portfolio or otherwise). The Report is made available with the understanding and expectation that each user will, with due care and diligence, conduct its own investigations and evaluations, and seek its own professional advice, in considering the Company's financial performance, strategies, prospects or risks, and the suitability of any investment therein for purchase, holding or sale within their portfolio.}
		
		\textbf{IEA Scenarios:}
		The choice of any scenario should not be taken as any endorsement of it, nor any statement as to the accuracy or completeness of its methodologies or assumptions, nor as a general preference for it over any other economic scenario. 2°ii may carry out the CCSA using other economic scenarios. Users must form their own view as to the decarbonisation or economic scenarios, trajectories and models that are most appropriate to their circumstances. No explicit or implicit assumption is made in relation to the current or future alignment of the scenarios with the Paris Agreement or the climate-related policies of any government at international, national or sub-national level.
		
		\textbf{TCFD and other climate-related financial reporting:}
		The CCSA and this Report may provide underlying data for climate-related financial reporting, including initiatives undertaken with regard to the Recommendations of G20 Financial Stability Board’s Taskforce on Climate-Related Financial Disclosures (TCFD). However, its use in isolation does not purport to provide ‘TCFD compliance’ or compliance with any other reporting requirements, and any data used should be subject to appropriate processes of verification and assurance. 
		
		\textbf{Exclusion of liability:}
		To the extent permitted by law we will not be liable to any user or to the Company for any direct, indirect or consequential loss or damage, whether in contract, tort (including negligence), breach of statutory duty or otherwise, even if foreseeable, relating to any information, data, content or opinions stated in this Report, or arising under or in connection with the use of, or reliance on, the CCSA or this Report.
		
		
		
	}	
	
	\PageFooter{Legal Disclaimer}
	\newpage
	%%%%%%%%%% 2 %%%%%%%%%%
	\section*{} % Executive Summary
	\HeaderDouble{Executive Summary}{\textsl{Company.Name}}
	
	This report by 2° Investing Initiative provides an assessment of \(\Companyname\)'s automotive production by technology, its future alignment with climate transition pathways and evaluates its performance against other automotive companies globally.
	
	\vspace{0.5cm}
	\Subheader{production in Start.Year} \\
	\(\Companyname\) is predicted to build CapInstalledStart vehicles with CompPercLowCStart\% coming from low carbon technologies - electric and hybrid. By End.Year, PlannedAdditions
	
	
	%Table 1
	\vspace{0.3cm} 
	
	\begin{center} % Table1 
		\adjincludegraphics[width = 1\linewidth,trim={0cm 0cm 0cm 0cm},clip]{CompanyReportOutputs/Table1.png} 
	\end{center}
	
	
	% Bar chart 
	\vspace{0.5cm}
	\Subheader{Comparison of production plans with transition scenarios}
	
	\(\Companyname\)’s planned additional production is compared to the different climate scenarios of the International Energy Agency (IEA) for each technology as described on page 4. The additional production planned by \(\Companyname\) ScenarioStatement
	
	\vspace{0.7cm}
	\begin{center} %bar chart
		\adjincludegraphics[width = 0.7\linewidth,trim={0cm 0cm 0cm 0cm},clip]{CompanyReportOutputs/Fig31.png} 
	\end{center}
	
	%Table 2
	\vspace{1cm}
	\Subheader{Changes in production required to align with the Select.Scenario by End.Year}\\
	In order to align with the ScenarioFullName (Select.Scenario) by End.Year, \(\Companyname\) would require the following additional changes in production by technology to its current plans by End.Year:
	\vspace{0.3cm}
	
	
	\begin{center} % Table1 
		\adjincludegraphics[width = 1\linewidth,trim={0cm 0cm 0cm 0cm},clip]{CompanyReportOutputs/Table2.png} 
	\end{center}
	
	\PageFooter{Executive Summary}
	\newpage
	
	
	
	%%%%%%%%%% 3 %%%%%%%%%%
	\section*{} % Introduction
	\HeaderSingle{Introduction}
	
	\Subheader{Key Questions}
	
	This climate scenario report addresses five key questions regarding \(\Companyname\)’s climate strategy:
	
	\begin{enumerate}[nolistsep]
		\item{How does the company’s current production mix compare to the automotive market’s production mix? (Page 6)}
		\item{How does the company’s future vehicle production compare to different climate transition scenarios? (Page 7)}
		\item{How does the company’s planned production mix by End.Year compare to the scenario-aligned market? 
			(Page 9)}
		\item{How can the company adjust its production plans to align with the Select.Scenario by End.Year? 
			(Page 10)}
		\item{How does the company’s climate alignment compare to other automotive companies? 
			(Page 11)}
	\end{enumerate} 
	
	This document solely presents the results of the above analyses. For more information on the methodology, sce\-nar\-i\-os, underlying data, and limitations, please refer to methodology documentation available at: \\ \href{https://www.transitionmonitor.com/company-reports/}{www.transitionmonitor.com/company-reports/}.
	
	\Subheader{Why is scenario analysis important?}\\
	Scenario analysis is highlighted within by the Task Force for Climate-related Financial Disclosures (TCFD) as a recommended tool for understanding the resilience of organization’s strategies under different climate related scenarios. It supports both companies and investors in developing action plans as a response to the Paris Agreement.
	
	\Subheader{How does this scenario analysis work?}\\
	This scenario analysis is an assessment of the physical assets owned by  \(\Companyname\)  and its plans for future production, based on third party data. The share of responsibility, defined by climate scenarios that outline possible transition pathways, has been allocated to the company according to the regional distribution of its automotive production. Further analyses allow us to understand how  \(\Companyname\) is currently, and in the future, exposed to climate transition risks and opportunities.
	
	\Subheader{How can it be used?}\\
	For \textbf{Companies}, this analysis provides a comparison of its performance relative to peers, and an understanding of how climate change responses differ. It also provides an overview of how planned production changes compare to the climate scenarios developed by the International Energy Agency (IEA). It highlights potential areas for action by companies.
	
	For \textbf{Investors}, this report may be used to inform their decision making by highlighting the alignment of the trajectories of companies in their portfolio with different climate scenarios and therefore their potential exposure to transition risks. The information provided in this report can support engagement activities with companies and may provide data for reporting requirements.
	
	For \textbf{other stakeholders}, such as policy makers or NGOs, this may support the development of guidelines for reporting or research.
	
	\textbf{What this report doesn’t do:} this report is not a financial analysis of the company and should not be taken as investment advice.
	
	Data used in this report is based on third party data from Auto Forecast Solutions (effective as of 06/2019) and may vary from what is announced by the company in annual reports; the data in this report	reflects an aggregation of the known subsidiaries of \(\Companyname\) aggregated under the equity share	principle. Details regarding the data sources and processing can be found on page 14. Companies are invited to review the data and provide feedback to assist in improving the underlying data sets by emailing 2dii at transitionmonitor@2degrees-investing.org. 
	
	
	\PageFooter{Introduction}
	\newpage
	
	
	
	%%%%%%%%%% 4 %%%%%%%%%%
	\section*{} % Reading the Report
	\HeaderSingle{Reading the Report}
	
	\Subheader{Report Contents}\\	
	This report consists of three elements:
	\vspace{-0.2cm}
	\begin{enumerate}[itemsep=0pt,parsep=2pt,topsep=2pt]
		\item{\textbf{Company profile:} information about the current installed production of the company, its technology mix and its global production distribution.} 
		\item{\textbf{Scenario Analysis:} results of the comparison of the company production plans to different scenarios and the market.}
		\item{\textbf{Peer Comparison:} a comparison of the scenario analysis results to peer companies operating in the same market.}
	\end{enumerate}
	%\par
	\vspace{-0.15cm}
	\Subheader{Key Concepts}\\
	To understand the results presented in this report, some of the key concepts are summarised below. For detailed
	information about the methodology, scenarios and underlying data, please refer to methodology documentation available at: \par \href{https://www.transitionmonitor.com/company-reports/}{www.transitionmonitor.com/company-reports/}.\par
	
	%\vspace{0.2cm}
	\textbf{\emph{Low carbon technologies:}} This analysis treats electric and hybrid vehicles as low carbon technologies, and internal combustion engine vehicles (ICE) as a high carbon technology. 
	
	\textbf{\emph{Hybrid Vehicles:}} The hybrid category considers both plug in hydrids and standalone vehicles.  
	
	%\vspace{0.2cm}
	\textbf{\emph{production:}} Refers to the production of new vehicles based upon a model developed by AFS. 
	
	%\vspace{0.2cm}
	\textbf{\emph{Production mix:}} The distribution of the automotive production of \(\Companyname\) is used as an indicator. This refers to the share of installed production that \(\Companyname\) has in each technology.\par
	
	%\vspace{0.2cm}
	\textbf{\emph{Market:}} The market referred to in this report is global and therefore includes all automotive companies globally.
	
	%\vspace{0.2cm}
	\textbf{\emph{Aligned with a scenario:}} To be aligned with a scenario implies that the future production of the company matches what is expected based on the roadmaps developed by the IEA.\par
	
	%\vspace{0.2cm}
	\textbf{\emph{Scenarios:}} Three IEA scenarios are included in this report's analysis: these are sourced from the Energy Technology Perspectives 2017 (ETP 2017) and are detailed in Table 1. These have been chosen due to their regional and technological gra\-nu\-la\-ri\-ty. The Select.Scenario is used as the benchmark scenario. The scenarios consist of technology roadmaps that outline the production changes required for each technology globally.This is aggregated to the company to determine the overall expected change required.\par
	
	
	\vspace{0.3cm}
	\Tab{1:} \TabCap{Overview of the IEA scenarios used in the analysis.}
	
	\begin{center} % Table1 
		\adjincludegraphics[width = 0.8\linewidth,trim={0cm 0cm 0cm 0cm},clip]{CompanyReportOutputs/TableScenarios.png} 
	\end{center}
	\vspace{0.1cm}
	\noindent\rule{8cm}{0.2pt}\\
	
	\Footnotetext{*The temperature rise estimates for the B2DS, and 2DS are specified by the IEA. The RTS estimate is taken from Climate Action Tracker's 2018 Warming Projections Global Update.}
	
	\PageFooter{Reading the Report}
	\newpage
	
	
	
	%%%%%%%%%% 5 %%%%%%%%%%
	\section*{} % Company Profile
	\HeaderSingle{Company Profile}
	
	This section outlines the current and future production mix of \(\Companyname\). Figure 1.1 shows the changes in production in each technology between Start.Year and End.Year. From this, one may be able to extrapolate whether the company’s transition risks increase or decrease. Figure 1.2 and 1.3 show the geographical distribution of automotive production and production mix.\par
	
	%\noindent\rule{8cm}{0.2pt}
	\vspace{-0.3cm}
	\begin{center}
		\hspace*{1.25cm}\includegraphics[width = 1\linewidth,trim={0cm 0cm 0cm 0cm},clip]{CompanyReportOutputs/Fig11.png}
	\end{center}
	\Fig{1.1:} \FigCap{Company production mix in Start.Year and End.Year.}
	
	%\noindent\rule{8cm}{0.2pt}
	\vspace{0.2cm}	
	\begin{center}
		\adjincludegraphics[width = 0.85\linewidth,trim={0cm 0cm 0cm 2cm},clip]{CompanyReportOutputs/Fig12.png}
	\end{center}
	\Fig{1.2:} \FigCap{Geographical distribution of the company’s automotive production in Start.Year.}
	
	%\noindent\rule{8cm}{0.2pt}
	\vspace{-0.2cm}	
	\begin{center}
		\hspace*{-0.15cm}\includegraphics[width = 1\linewidth,trim={0cm 0cm 0cm 0cm},clip]{CompanyReportOutputs/Fig13.png}
	\end{center}
	\Fig{1.3:} \FigCap{Overview of the company’s production mix and total production in the largest countries by total production in Start.Year.}	
	
	\PageFooter{Company Profile}
	\newpage
	
	
	%%%%%%%%%%6 %%%%%%%%%%
	\section*{} % Current Alignment
	\HeaderSingle{Current Alignment}
	
	\Subheader{How does the current production mix of  \(\Companyname\) compare to the market?}\\
	This section provides an overview of the diversification of  \(\Companyname\)'s production across high and low carbon te\-chno\-lo\-gies. In order to meet the goals of the Paris Agreement, the IEA broadly signals that the share of “low carbon technologies” must increase while the share of “high carbon technologies” must decrease.
	
	The company's production mix is presented in terms of the breakdown of its production by technology. The market is representative of all automotive companies in the global automotive market. 
	
	
	\vspace{0.5cm}	
	%\begin{center}
	\hspace*{0.5cm}\textbf{Low and high carbon production mix percentage}
	
	\hspace*{1.5cm}\includegraphics[width = 0.7\linewidth,trim={0cm 0cm 0cm 1cm},clip]{CompanyReportOutputs/Fig4_Start.Year.png}\par
	%\end{center}
	
	
	\Fig{2.1:} \FigCap{Comparison of the company's production mix to the market's production mix in Start.Year by technology.}	
	
	\vspace{0.2cm}	
	\ColortextLarge{\(\Companyname\) has CompPercLowCStart\% of its automotive production in low carbon technologies compared to MarketPercLowCStart\% in the market. \(\Companyname\) TechShareComparisonStart}
	
	\PageFooter{Current Alignment}
	\newpage
	
	
	
	%%%%%%%%%% 7 %%%%%%%%%%
	\section*{} % Trajectory
	\HeaderSingle{Trajectory}
	
	\Subheader{How do the capital expenditure plans for different technologies compare to the climate sce\-narios?}\par
	\(\Companyname\)'s plans to increase or decrease production in different technologies over the next 5 years can be compared to the In\-ter\-na\-tio\-nal Energy Agency (IEA) scenarios. These scenarios present possible transition pathways and the
	changes in production required if each company in the world were to align its production accordingly.\par
	The expected change in production by technology as per the IEA scenarios has been applied to the automotive production of
	\(\Companyname\) to calculate the changes required under each scenario. This report benchmarks the
	company against the ScenarioFullName (Select.Scenario), though the following charts show the Beyond 2 Degree Scenario (B2DS), the 2 Degree Scenario (2DS) and the Reference Technology Scenario (RTS) together.\par
	Alignment with climate scenarios may vary by technology. For each technology, figure 3.1 summarises the different IEA sce\-na\-rios
	that \(\Companyname\)’s future production plan aligns with. It is important to note that these charts are independent of the current exposure to each technology
	(that only determines the starting point in terms of production). The initial (Start.Year) weighting of a technology within
	the company’s production mix is not reflected in these charts.
	\vspace{0.3cm}
	\begin{center} % tile chart
		\adjincludegraphics[width = 0.7\linewidth,trim={0cm 0cm 0cm 0cm},clip]{CompanyReportOutputs/Fig31.png} 
	\end{center}
	\vspace{0.3cm}
	\Fig{3.1:} \FigCap{Scenario outcome of the change in production plans for each technology by End.Year. This summarises the results of the trajectory charts in End.Year.}
	
	\vspace{0.3cm}
	\ColortextLarge{The additional production planned by \(\Companyname\) ScenarioStatement}\par
	The charts on the following page (figure 3.2) provide additional details on how \(\Companyname\)’s
	production plans for each technology align with three IEA scenarios over the next five years. They also show the
	market’s trajectory for context.
	
	\vspace{0.5cm}
	%\noindent\rule{2.5cm}{0.2pt}
	\begin{multicols}{2}
		%\Subheader{Annotation}\\
		\textsl{The background colours represent how the trajectory of a technology should progress under the relevant scenario based on the company's current automotive production. The solid and dashed lines represent the production plans of the company and those of the automotive market scaled to the starting point of the company. In the chart to the right, the company's production plans for this technology lie between the SDS and NPS trajectories. The difference in End.Year between the company's production plan and the end point for a specific scenario indicates the change in production that would be required for alignment. The market production can be compared to the company plans as a relative indicator only, as the scenarios are specific to the company. In this case, the company will produce relatively more of this vehicle type than the market.}\par
		
		
		\begin{center} % trajectory gas as example
			\adjincludegraphics[width = 0.5\textwidth,trim={0cm 0cm 0cm 0.7cm},clip]{CompanyReportGraphics/246Explainer2.png} 
			
		\end{center}
	\end{multicols}
	
	\PageFooter{Trajectory}
	\newpage
	
	
	
	%%%%%%%%%% 8 %%%%%%%%%%
	\section*{} % Trajectoery
	\HeaderSingle{Trajectory}
	
	\vspace{-0.2cm}
	\begin{center} %trajectory 
		\adjincludegraphics[width = .5\linewidth,trim={0cm 0cm 0cm 0.5cm},clip]{CompanyReportOutputs/Trajectory_Electric.png} 
		\adjincludegraphics[width = .5\linewidth,trim={0cm 0cm 0cm 0.5cm},clip]{CompanyReportOutputs/Trajectory_Hybrid.png} 
		\adjincludegraphics[width = .5\linewidth,trim={0cm 0cm 0cm 0.5cm},clip]{CompanyReportOutputs/Trajectory_ICE.png} 
	\end{center}
	
	
	
	
	\vspace{-0.2cm}
	\begin{center} %Trajectory Legend
		\adjincludegraphics[width = 0.6\linewidth,trim={0.5cm 0.2cm 0.2cm 0.5cm},clip]{CompanyReportGraphics/TrajectoryLegend_ETP.png} 
	\end{center}
	% \vspace{0.5cm}
	\Fig{3.2: }\FigCap{illustrates how \(\Companyname\)’s planned production changes in each technology compare to different IEA transition pathways and the market.}
	
	\PageFooter{Trajectory}
	\newpage
	
	
	
	%%%%%%%%%% 9 %%%%%%%%%%
	\section*{} % Future Alignment
	\HeaderSingle{Future Alignment}
	
	\Subheader{How will the planned production mix of \(\Companyname\) compare to a global market aligned with the Select.Scenario in End.Year?}\par
	The production mix of \(\Companyname\)  in End.Year is based on its production mix in Start.Year plus planned production changes between Start.Year and End.Year. The aligned market production mix shows what would be expected if the current global automotive market were to develop over the next five years in accordance with the Select.Scenario.\par
	
	If the company has a lower amount of low carbon technologies than the theoretical aligned market, it may be exposed to higher transition risks based on the technological trajectories outlined by the IEA.\par
	
	Figure 4 shows that \(\Companyname\) has an production mix in End.Year which has FutureLowHigh\% percentage points or RelativeHighLow\% LessOrMore low carbon production than an aligned market.
	
	\vspace{0.5cm}	
	%\begin{center}
	\hspace*{0.5cm}\textbf{Low and high carbon production mix percentage}
	
	\hspace*{1.5cm}\includegraphics[width = 0.7\linewidth,trim={0cm 0cm 0cm 0cm},clip]{CompanyReportOutputs/Fig2_End.Year.png}\par
	%\end{center}
	%\begin{center}
		
	\Fig{4.1:} \FigCap{Comparison of the company production mix to the market production mix in End.Year.}	
	
	\vspace{0.5cm}	
	\ColortextLarge{By End.Year \(\Companyname\) TechShareComparisonEnd} \par 
	
	\PageFooter{Future Alignment}
	\newpage
	
	
	
	
	%%%%%%%%%% 10 %%%%%%%%%%	
	\section*{} % Achieving Alignment
	
	\HeaderSingle{Achieving Alignment}
	
	\Subheader{What changes in production are required by \(\Companyname\) to align itself with the Select.Scenario?}
	
	For \(\Companyname\) to align itself with the Select.Scenario by End.Year based on the company’s current production, the following production changes by technology are required. 
	
	\vspace{0.5cm}	
	\begin{center}
		\hspace*{0.4cm}\includegraphics[width = 1.2\linewidth,trim={0cm 0cm 0cm 0cm},clip]{CompanyReportOutputs/Fig5.png}
		\adjincludegraphics[width = 0.8\linewidth,trim={0cm 0cm 0cm 0cm},clip]{CompanyReportGraphics/LegendFig5.png} 
	\end{center}
	\Fig{5:} \FigCap{Changes in production required to align with the Select.Scenario.}	
	
	\vspace{0.5cm}	
	\ColortextLarge{By End.Year, \(\Companyname\) AchievingAlignment}\par
	In some cases, the company’s production plans may outperform the production required to align with the Select.Scenario. If the company’s production plan for low carbon technologies exceeds scenario targets, no changes in production plans are specified. Similarly, no additional production is specified if the company’s plans already meet the transition pathways for ICE vehicles.
	
	\PageFooter{Achieving Alignment
	}
	\newpage
	
	
	
	
	%%%%%%%%%% 11 %%%%%%%%%%
	\section*{} % Comparison Between Automotive Companies
	\HeaderSingle{Comparison Between automotive companies}
	
	\Subheader{How does the current production and future planned production of \(\Companyname\) for low carbon tech\-no\-lo\-gies compare to other automotive companies in the global market?}\par
	In this section, we represent the current and future production mix of \(\Companyname\) relative to the other automotive companies in the global market. Figure 6 highlights: 
	%\begin{enumerate}
	\begin{itemize}
		\item On the x-axis, the percentage of low carbon technologies in the production mix in Start.Year.
		\item On the y-axis, the percentage of low carbon technologies in the production mix in End.Year.
		\item The Start.Year total automotive production of each company via the size of the circles. Each circle represents a separate automotive company.
	\end{itemize}
	%\end{enumerate}
	\vspace{0.2cm}
	\(\Companyname\) is highlighted in black.
	\vspace{-0.2cm}
	\begin{center} %Trajectory Legend
		\adjincludegraphics[width = 0.9\linewidth,trim={0cm -1cm 0cm 2cm},clip]{CompanyReportOutputs/Fig6.png} 
	\end{center}
	
	\Fig{6:} \FigCap{\% planned future production in low carbon technologies vs \% current production in low carbon technologies.}	
	
	\vspace{0cm}	
	When the company lies above the diagonal line it indicates that the share of low carbon technology is increasing. The further up the company is, the higher their low carbon share is in the future. The further to the right the company is, the higher is their current share of low carbon production.  
	
	
	\PageFooter{Comparison Between automotive companies}
	\newpage
	
	
	
	
	%%%%%%%%%% 12 %%%%%%%%%%
	

	
	%%%%%%%%%% 13 %%%%%%%%%%
	\section*{} % Market Share
	
	\HeaderSingle{Market Share}
	
	\Subheader{How does \(\Companyname\)’s market share of each automotive technology evolve between Start.Year and End.Year?
	}\\
	
	This section shows how \(\Companyname\)'s market share in low carbon technologies is expected to develop between Start.Year and End.Year and what the company's future positioning relative to the market will be. \par
	
	Figure 8 shows changes in the company’s electric and total automotive production market share, defined as the \% of total production in each technology over the entire global automotive production. 
	
	\vspace{0cm}
	\begin{center} % comparison utlities bar chart 
		\adjincludegraphics[width = 1\linewidth,trim={0.5cm 0cm 0cm 0.5cm},clip]{CompanyReportOutputs/Fig8.png} 
	\end{center}
	
	\vspace{0cm}
	\Fig{8:} \FigCap{The above charts show how \(\Companyname\)’s share of the global market develops for electric and total automotive production. The chart on the left provides a comparison to the global market given current plans, and the chart on the right shows how this would develop if the market were aligned with the Select.Scenario.}	
	
	\vspace{0.5cm}	
	\ColortextLarge{The total market share of \(\Companyname\) ChangeMS Its market share of electric vehicle production RenewablesMSChange}
	
	The market share for each technology represents the company’s production as a percentage of the production of all automotive companies (actual and aligned) including current announced plans. If the company’s electric vehicle market share is decreasing over the next 5 years, this suggests that \(\Companyname\) plans to increase electric vehicle production at a lower rate than the automotive companies market as a whole.
	
	
	\PageFooter{Market Share}
	\newpage
	
	
	
	
	%%%%%%%%%% 14 %%%%%%%%%%
	\section*{} % Data Sources
	
	\HeaderSingle{Data Sources}
	
	2dii scenario analysis is based on forward-looking automotive data from Auto Forecast Solutions (AFS) paired with company ownership information from Bloomberg. 
	
	AFS use a model to predict future production values for plants and therefore countries at a company level. AFS provides granular information on the brand, nameplate, drive train technology, and location of production currently and modelled for the upcoming eight years.  This modelling takes into consideration not just what the company states online and publically but also considers geopolitical, consumer behaviour, competition and shareholder values. The modelling generally takes a conservative approach. Backtesting of this model, has shown the model production estimates have a varience of up to 5\% of total global production with a 5 year time frame. Nevertheless there may be discrepancies at the brand nameplate level. 
	
	The company’s current production mix was calculated by aggregating production where the company is listed as owner, weighted by an ownership stake. 5-year production plans were calculated by aggregating production for the years between Start.Year and End.Year. The result is a forward-looking production mix for \(\Companyname\) that serves as starting point and basis for comparison for scenario analysis. 
	
	
	Reasons for variation from company reported production may fall into the following categories:
	
	\begin{enumerate}
		\item We take asset data and ownership information from two major data providers: AFS and Bloomberg. While the production data at a global level has a low error margin, this may be more for any given company or technology. 
		\item  We allocate production from subsidiaries to parent companies according to the following rules: If a subsidiary company is private/unlisted, 100\% of its production is allocated to the parent company holding the controlling stake. If a subsidiary is public/listed, the non-free float portion of its production is allocated to the parent company holding the controlling stake. No automotive production is allocated to parent companies holding non-controlling stakes.
	\end{enumerate}
	
	Please review the legal disclaimer for further information about the limitations of the data.
	
	
	\PageFooter{Data Sources}
	\newpage
	
	
	
	
	
	
\end{document} 






